% (c) Nikita Lisitsa, lisyarus@gmail.com, 2026

\documentclass[10pt]{beamer}

\usepackage[T2A]{fontenc}
\usepackage[russian]{babel}
\usepackage{minted}

\usepackage[most]{tcolorbox}
\tcbuselibrary{minted}

\usepackage{graphicx}
\graphicspath{ {./images/} }

\usepackage{adjustbox}

\usepackage{color}
\usepackage{soul}
\usepackage{multicol}
\usepackage{hyperref}
\usepackage{tikz}
\usetikzlibrary{arrows.meta,positioning,calc}

\usetheme{metropolis}

\definecolor{red}{rgb}{1,0,0}
\definecolor{codebg}{RGB}{29,35,49}
\setminted{fontsize=\footnotesize}

\makeatletter
\newcommand{\slideimage}[1]{
  \begin{figure}
    \begin{adjustbox}{width=\textwidth, totalheight=\textheight-2\baselineskip-2\baselineskip,keepaspectratio}
      \includegraphics{#1}
    \end{adjustbox}
  \end{figure}
}
\makeatother

\title{Язык программирования C++}
\subtitle{Лекция 20: Race condition, mutex, atomic, регулярные выражения, static\_assert, compile-time вычисления, variadic templates}
\date{2026}
\author{Никита Лисица}

\setbeamertemplate{footline}[frame number]

\begin{document}

\frame{\titlepage}

\begin{frame}[fragile]
\frametitle{Race condition}
\begin{itemize}
\usemintedstyle{tango}
\item С точки зрения потока/процесса, переключение контекста может произойти в произвольный момент времени (в том числе в середине выполнения какой-нибудь операции)
\pause
\item Как правило, мы не можем явно управлять распределением потоков по ядрам CPU, порядком их выполнения, и т.д.
\pause
\item Если результат выполнения программы зависит от порядка, в котором выполнялись инструкции в разных потоках, говорят о \textit{состоянии гонки (race condition)}
\end{itemize}
\end{frame}

\begin{frame}[fragile]
\frametitle{Race condition: пример}
\begin{itemize}
\usemintedstyle{lightbulb}
\begin{tcblisting}{listing engine=minted, minted language=cpp, minted options={fontsize=\scriptsize}, listing only, colback=codebg, colframe=codebg, rounded corners}
#include <thread>
#include <iostream>

void hello() {
  std::cout << "Hello from " << std::this_thread::get_id()
    << std::endl;
}

int main() {
  std::vector<std::thread> threads;
  for (int i = 0; i < 5; ++i)
    threads.emplace_back(&hello);
  for (auto & thread : threads)
    thread.join();
}
\end{tcblisting}
\pause
\item Что выведется на экран?
\end{itemize}
\end{frame}

\begin{frame}[fragile]
\frametitle{Race condition: пример}
\begin{itemize}
\item Может быть
\usemintedstyle{lightbulb}
\begin{tcblisting}{listing engine=minted, minted language=cpp, minted options={fontsize=\scriptsize}, listing only, colback=codebg, colframe=codebg, rounded corners}
Hello from 140276650997504
Hello from 140276667782912
Hello from 140276659390208
Hello from 140276642604800
Hello from 140276676175616
\end{tcblisting}
\pause
\item А может быть
\usemintedstyle{lightbulb}
\begin{tcblisting}{listing engine=minted, minted language=cpp, minted options={fontsize=\scriptsize}, listing only, colback=codebg, colframe=codebg, rounded corners}
Hello from Hello from Hello from 
140276650997504Hello
from 140276659390208Hello from 
140276667782912
140276676175616
140276642604800
\end{tcblisting}
\end{itemize}
\end{frame}

\begin{frame}[fragile]
\frametitle{Race condition: пример}
\begin{itemize}
\item Один поток добавляет элементы в связный список, второй поток итерируется по нему
\pause
\usemintedstyle{lightbulb}
\begin{tcblisting}{listing engine=minted, minted language=cpp, minted options={fontsize=\scriptsize}, listing only, colback=codebg, colframe=codebg, rounded corners}
auto next = curr.next;
curr.next = new_node;
// Поток #1 остановился здесь
new_node.next = next;
\end{tcblisting}
\pause
\usemintedstyle{tango}
\item Во втором потоке сломается итерация по списку (например, если \mintinline{cpp}|new_node.next| было не проинициализировано)
\end{itemize}
\end{frame}

\begin{frame}[fragile]
\frametitle{Race condition: пример}
\begin{itemize}
\usemintedstyle{lightbulb}
\begin{tcblisting}{listing engine=minted, minted language=cpp, minted options={fontsize=\scriptsize}, listing only, colback=codebg, colframe=codebg, rounded corners}
int x = 0;

// Поток #1
while (!stopped) {
  ++x;
}

// Поток #2
while (!stopped) {
  if (x % 2 == 0) {
    std::cout << x;
  }
}
\end{tcblisting}
\pause
\usemintedstyle{tango}
\item Может быть выведено нечётное число
\pause
\item Поток \#2 выполнил проверку и зашёл внутрь \mintinline{cpp}|if|, затем передал управление потоку \#1
\pause
\item Поток \#1 изменил \mintinline{cpp}|x| на нечётное значение и передал управление потоку \#2
\end{itemize}
\end{frame}

\begin{frame}[fragile]
\frametitle{Гонка данных (data race)}
\begin{itemize}
\usemintedstyle{tango}
\item Когда один поток пишет в некоторую память, и одновременно один или несколько потоков эту память читают, говорят о \textit{гонке данных (data race)}
\pause
\item В общем случае результат гонки данных -- неопределённое поведение
\pause
\item На некоторых архитектурах некоторые инструкции атомарны (напр. чтение или запись до 8 байт по выровненному адресу на x86-64)
\pause
\item Как правило, операции в языке программирования компилируются в \textit{несколько} инструкций CPU (\(\Longrightarrow\) не являются атомарными)
\pause
\item Оптимизации компилятора могут перемешивать инструкции, меняя поведение в многопоточном случае
\end{itemize}
\end{frame}

\begin{frame}[fragile]
\frametitle{Гонка данных: пример}
\begin{itemize}
\usemintedstyle{lightbulb}
\begin{tcblisting}{listing engine=minted, minted language=cpp, minted options={fontsize=\scriptsize}, listing only, colback=codebg, colframe=codebg, rounded corners}
int x = 0;

std::vector<std::thread> threads;
for (int i = 0; i < 5; ++i)
  threads.emplace_back([&x]{
    for (int i = 0; i < 1000; ++i)
      ++x;
  });

for (auto & thread : threads)
  thread.join();

std::cout << x << std::endl;
\end{tcblisting}
\pause
\item Что выведется на экран?
\pause
\item В общем случае -- что угодно (UB), на x86-64 -- что угодно от 1000 до 5000
\end{itemize}
\end{frame}

\begin{frame}[fragile]
\frametitle{Гонка данных: пример}
\begin{itemize}
\usemintedstyle{tango}
\item \mintinline{cpp}|++x| это, на самом деле, три операции:
\pause
\begin{itemize}
\item Прочитать \mintinline{cpp}|x| из памяти в регистр
\pause
\item Инкрементировать регистр
\pause
\item Записать значение регистра обратно в память
\end{itemize}
\pause
\item Возможное исполнение:
\pause
\begin{itemize}
\item Поток \#1 прочитал \mintinline{cpp}|x|
\pause
\item Поток \#2 прочитал \mintinline{cpp}|x|
\pause
\item Оба потока инкрементировали регистр
\pause
\item Поток \#1 записал \mintinline{cpp}|x|
\pause
\item Поток \#2 записал \mintinline{cpp}|x|
\pause
\item \(\Longrightarrow\) \mintinline{cpp}|x| изменился с 0 на 1
\end{itemize}
\end{itemize}
\end{frame}

\begin{frame}[fragile]
\frametitle{Мьютекс}
\begin{itemize}
\usemintedstyle{tango}
\item \textit{Мьютекс (mutex, от англ. mutual exclusion)} -- примитив синхронизации, ограничивающий доступ потоков к данным или коду
\pause
\item Часть кода, защищённая мьютексом (т.н. \textit{критическая секция}) может одновременно выполняться \textbf{только одним потоком} (остальные будут \textbf{ждать} своей очереди)
\pause
\item В стандартной библиотеке есть реализация мьютекса \mintinline{cpp}|std::mutex|
\end{itemize}
\end{frame}

\begin{frame}[fragile]
\frametitle{Мьютекс}
\begin{itemize}
\usemintedstyle{lightbulb}
\begin{tcblisting}{listing engine=minted, minted language=cpp, minted options={fontsize=\scriptsize}, listing only, colback=codebg, colframe=codebg, rounded corners}
std::mutex mutex;
int x = 0;

// ...

threads.emplace_back([&x]{
  for (int i = 0; i < 1000; ++i) {
    // Код между lock и unlock - критическая секция
    mutex.lock();
    ++x;
    mutex.unlock();
  }
});
\end{tcblisting}
\end{itemize}
\end{frame}

\begin{frame}[fragile]
\frametitle{Мьютекс: lock\_guard}
\begin{itemize}
\usemintedstyle{tango}
\item Почти всегда \mintinline{cpp}|mutex.lock()| и \mintinline{cpp}|mutex.unlock()| идут в паре
\pause
\item Чтобы не забыть сделать \mintinline{cpp}|mutex.unlock()| (иначе потоки могут зависнуть!), есть \mintinline{cpp}|std::lock_guard|
\usemintedstyle{lightbulb}
\begin{tcblisting}{listing engine=minted, minted language=cpp, minted options={fontsize=\scriptsize}, listing only, colback=codebg, colframe=codebg, rounded corners}
// Сделает lock() в конструкторе и unlock() в деструкторе
std::lock_guard<std::mutex> lock{mutex};
++x;
\end{tcblisting}
\pause
\begin{tcblisting}{listing engine=minted, minted language=cpp, minted options={fontsize=\scriptsize}, listing only, colback=codebg, colframe=codebg, rounded corners}
// Можно воспользоваться CTAD
std::lock_guard lock{mutex};
++x;
\end{tcblisting}
\end{itemize}
\end{frame}

\begin{frame}[fragile]
\frametitle{Мьютекс: lock\_guard}
\begin{itemize}
\usemintedstyle{tango}
\item Можно отдать в \mintinline{cpp}|std::lock_guard| уже заблокированный мьютекс:
\usemintedstyle{lightbulb}
\begin{tcblisting}{listing engine=minted, minted language=cpp, minted options={fontsize=\scriptsize}, listing only, colback=codebg, colframe=codebg, rounded corners}
mutex.lock();
// Сделает unlock() в деструкторе
std::lock_guard lock{mutex, std::adapt_lock};
++x;
\end{tcblisting}
\end{itemize}
\end{frame}

\begin{frame}[fragile]
\frametitle{Мьютекс: unique\_lock}
\begin{itemize}
\usemintedstyle{tango}
\item Если может потребоваться сбросить состояние заблокированности до деструктора, или переместить (move) владение блокировкой, есть \mintinline{cpp}|std::unique_lock|
\pause
\usemintedstyle{lightbulb}
\begin{tcblisting}{listing engine=minted, minted language=cpp, minted options={fontsize=\scriptsize}, listing only, colback=codebg, colframe=codebg, rounded corners}
std::unique_lock lock{mutex};
++x;
if (x > 10) {
  // Можно вызвать unlock() самим
  lock.unlock();
  return;
}
++x;
\end{tcblisting}
\end{itemize}
\end{frame}

\begin{frame}[fragile]
\frametitle{Мьютекс: unique\_lock}
\begin{itemize}
\usemintedstyle{lightbulb}
\begin{tcblisting}{listing engine=minted, minted language=cpp, minted options={fontsize=\scriptsize}, listing only, colback=codebg, colframe=codebg, rounded corners}
std::unique_lock lock{mutex, std::defer_lock};
// После этого мьютекс будет заблокирован
lock.lock();
\end{tcblisting}
\pause
\usemintedstyle{lightbulb}
\begin{tcblisting}{listing engine=minted, minted language=cpp, minted options={fontsize=\scriptsize}, listing only, colback=codebg, colframe=codebg, rounded corners}
std::unique_lock lock{mutex, std::try_lock};
if (lock.owns_lock()) {
  // Получилось заблокировать, можно войти
  // в критическую секцию
} else {
  // Не получилось заблокировать - другой поток
  // уже находится в критической секции
}
\end{tcblisting}
\end{itemize}
\end{frame}

\begin{frame}[fragile]
\frametitle{Взаимная блокировка (deadlock)}
\begin{itemize}
\usemintedstyle{lightbulb}
\begin{multicols}{2}
\begin{tcblisting}{listing engine=minted, minted language=cpp, minted options={fontsize=\scriptsize}, listing only, colback=codebg, colframe=codebg, rounded corners}
// Поток #1
mutex1.lock();
mutex2.lock();
\end{tcblisting}
\columnbreak
\begin{tcblisting}{listing engine=minted, minted language=cpp, minted options={fontsize=\scriptsize}, listing only, colback=codebg, colframe=codebg, rounded corners}
// Поток #2
mutex2.lock();
mutex1.lock();
\end{tcblisting}
\end{multicols}
\pause
\usemintedstyle{tango}
\item Поток \#1 заблокировал \mintinline{cpp}|mutex1|, а поток \#2 заблокировал \mintinline{cpp}|mutex2|
\pause
\item Теперь поток \#1 ждёт разблокировки \mintinline{cpp}|mutex2|, а поток \#2 ждёт разблокировки \mintinline{cpp}|mutex1|
\pause
\item \(\Longrightarrow\) мьютексы никогда не будут разблокированы, потоки зависнут
\pause
\item Эта ситуация называется \textit{взаимной блокировкой (deadlock)}
\end{itemize}
\end{frame}

\begin{frame}[fragile]
\frametitle{Взаимная блокировка (deadlock)}
\begin{itemize}
\usemintedstyle{tango}
\item Существуют алгоритмы предотвращения взаимной блокировки \textit{(deadlock avoidance)}
\pause
\item В стандартной библиотеке есть функция \mintinline{cpp}|std::lock(mutex1, mutex2, ...)|, реализующая такой алгоритм
\pause
\item Есть RAII-обёртка \mintinline{cpp}|std::scoped_lock|, делающая то же самое (и разблокировку в деструкторе)
\end{itemize}
\end{frame}

\begin{frame}[fragile]
\frametitle{Предотвращение взаимной блокировки}
\begin{itemize}
\usemintedstyle{lightbulb}
\begin{tcblisting}{listing engine=minted, minted language=cpp, minted options={fontsize=\scriptsize}, listing only, colback=codebg, colframe=codebg, rounded corners}
std::lock(mutex1, mutex2);
std::lock_guard lock1(mutex1, std::adopt_lock);
std::lock_guard lock2(mutex2, std::adopt_lock);
// ...
\end{tcblisting}
\pause
\usemintedstyle{lightbulb}
\begin{tcblisting}{listing engine=minted, minted language=cpp, minted options={fontsize=\scriptsize}, listing only, colback=codebg, colframe=codebg, rounded corners}
std::unique_lock lock1(mutex1, std::defer_lock);
std::unique_lock lock2(mutex2, std::defer_lock);
std::lock(lock1, lock2);
// ...
\end{tcblisting}
\pause
\usemintedstyle{lightbulb}
\begin{tcblisting}{listing engine=minted, minted language=cpp, minted options={fontsize=\scriptsize}, listing only, colback=codebg, colframe=codebg, rounded corners}
std::scoped_lock lock(mutex1, mutex2);
// ...
\end{tcblisting}
\end{itemize}
\end{frame}

\begin{frame}[fragile]
\frametitle{Рекурсивный мьютекс}
\begin{itemize}
\usemintedstyle{tango}
\item Обычный мьютекс нельзя заблокировать дважды одним и тем же потоком (undefined behavior, в некоторых реализациях -- \mintinline{cpp}|std::system_error|)
\pause
\item Для такого существует рекурсивный/reentrant мьютекс:
\usemintedstyle{lightbulb}
\begin{tcblisting}{listing engine=minted, minted language=cpp, minted options={fontsize=\scriptsize}, listing only, colback=codebg, colframe=codebg, rounded corners}
std::recursive_mutex mutex;

void foo() {
  std::lock_guard lock(mutex);
  ...
  if (...) {
    foo();
  }
}
\end{tcblisting}
\end{itemize}
\end{frame}

\begin{frame}[fragile]
\frametitle{Атомарные типы (atomics)}
\begin{itemize}
\usemintedstyle{tango}
\item Атомарные операции -- операции, которые выполняются целиком, не прерываясь и не перемешиваясь с другими операциями (в т.ч. других потоков)
\pause
\item Даже если архитектура CPU (напр. x86-64) гарантирует атомарность чтения/записи в каких-то случаях, компилятор может перемешать эти инструкции с другими, что повлияет на логику программы
\pause
\item В стандартной библиотеке есть шаблон \mintinline{cpp}|std::atomic|, реализующий атомарные операции
\pause
\usemintedstyle{lightbulb}
\begin{tcblisting}{listing engine=minted, minted language=cpp, minted options={fontsize=\scriptsize}, listing only, colback=codebg, colframe=codebg, rounded corners}
std::atomic<int> count = 0;

// Из нескольких потоков
for (int i = 0; i < 1000; ++i)
  // Атомарно прибавить 1
  count.fetch_add(1);
\end{tcblisting}
\end{itemize}
\end{frame}

\begin{frame}[fragile]
\frametitle{Атомарные типы (atomics)}
\begin{itemize}
\usemintedstyle{tango}
\item В \mintinline{cpp}|std::atomic| можно завернуть любой примитивный тип, указатель, и простые структуры
\pause
\item \mintinline{cpp}|load()| -- атомарно прочитать значение (возвращает копию)
\pause
\item \mintinline{cpp}|store(value)| -- атомарно записать значение
\pause
\item \mintinline{cpp}|exchange(value)| -- атомарно записать новое значение и вернуть старое
\pause
\item \mintinline{cpp}|compare_exchange_weak/strong(value)| -- реализация compare-and-swap (CAS)
\pause
\item Есть особые методы для числовых типов: \mintinline{cpp}|fetch_add(value)|, \mintinline{cpp}|fetch_sub(value)|, \mintinline{cpp}|fetch_min(value)|, \mintinline{cpp}|fetch_max(value)|, \mintinline{cpp}|fetch_and(value)|, \mintinline{cpp}|fetch_or(value)|, \mintinline{cpp}|fetch_xor(value)|
\end{itemize}
\end{frame}

\begin{frame}[fragile]
\frametitle{Атомарные типы (atomics)}
\begin{itemize}
\usemintedstyle{tango}
\item \mintinline{cpp}|std::atomic| может быть реализован через мьютекс, а может -- через специфичные для платформы механизмы
\pause
\item Проверить, что \mintinline{cpp}|atomic| не использует мьютексы: \mintinline{cpp}|is_lock_free()|, \mintinline{cpp}|is_always_lock_free|
\pause
\item \textbf{NB:} Мьютексы можно реализовать через \mintinline{cpp}|atomic| используя CAS (т.н. \textit{spinlock})
\end{itemize}
\end{frame}

\begin{frame}[fragile]
\frametitle{Мьютексы, atomics, и оптимизации}
\begin{itemize}
\usemintedstyle{tango}
\item Мьютексы/атомики не только исполняют специфичные инструкции процессора, но и говорят компилятору о том, что определённые операции нельзя перемешивать
\pause
\item Почти все функции работы с атомиками принимают аргумент \textit{memory order}, описывающий, какие именно гарантии нам нужны от этой операции (по умолчанию -- \textit{sequential consistency}, самая сильная гарантия)
\pause
\item Есть древний совет использовать \mintinline{cpp}|volatile| переменные в качестве атомиков -- в современных CPU И компиляторах это \textbf{не работает}!
\end{itemize}
\end{frame}

\begin{frame}[fragile]
\frametitle{Регулярные выражения (regex)}
\begin{itemize}
\usemintedstyle{tango}
\item В стандартной библиотеке есть реализация регулярных выражений -- класс \mintinline{cpp}|std::regex|
\pause
\item Поддерживает много вариаций регулярных выражений (ECMAScript, basic POSIX, extended POSIX, ...)
\pause
\usemintedstyle{lightbulb}
\begin{tcblisting}{listing engine=minted, minted language=cpp, minted options={fontsize=\scriptsize}, listing only, colback=codebg, colframe=codebg, rounded corners}
std::regex number_regex("(\d)+");
std::smatch match;
std::regex_search("1234", match, number_regex);
\end{tcblisting}
\pause
\usemintedstyle{tango}
\item \mintinline{cpp}|std::regex_match| -- проверка на соответствие регулярному выражению
\pause
\item \mintinline{cpp}|std::regex_search| -- поиск подстроки по регулярному выражению
\pause
\item \mintinline{cpp}|std::regex_replace| -- замена по регулярному выражению
\end{itemize}
\end{frame}

\begin{frame}[fragile]
\frametitle{static\_assert}
\begin{itemize}
\usemintedstyle{tango}
\item Позволяет сделать compile-time проверку некоторого условия:
\usemintedstyle{lightbulb}
\begin{tcblisting}{listing engine=minted, minted language=cpp, minted options={fontsize=\scriptsize}, listing only, colback=codebg, colframe=codebg, rounded corners}
// Текст ошибки опционален
static_assert(sizeof(int) == 4, "unsupported platform");

struct mesh_vertex {
  vec3f position;
  vec3f normal;
  vec2f texcoord;
};

static_assert(sizeof(mesh_vertex) == 32);

static_assert(std::atomic<int>::is_always_lock_free);
\end{tcblisting}
\end{itemize}
\end{frame}

\begin{frame}[fragile]
\frametitle{Compile-time вычисления}
\begin{itemize}
\usemintedstyle{tango}
\item Шаблоны в C++ можно рассматривать как функциональный язык программирования, выполняющийся в процессе компиляции (compile-time)
\pause
\usemintedstyle{lightbulb}
\begin{tcblisting}{listing engine=minted, minted language=cpp, minted options={fontsize=\scriptsize}, listing only, colback=codebg, colframe=codebg, rounded corners}
template <unsigned int N>
struct factorial {
  static const auto value = N * factorial<N - 1>::value;
};

template <>
struct factorial<0> {
  static const unsigned int value = 1;
};

auto x = factorial<5>::value; // x = 120
\end{tcblisting}
\end{itemize}
\end{frame}

\begin{frame}[fragile]
\frametitle{Compile-time вычисления}
\begin{itemize}
\usemintedstyle{tango}
\item Односвязный список типов:
\usemintedstyle{lightbulb}
\begin{tcblisting}{listing engine=minted, minted language=cpp, minted options={fontsize=\scriptsize}, listing only, colback=codebg, colframe=codebg, rounded corners}
struct empty{};

template <typename Head, typename Tail>
struct append {
  using head = Head;
  using tail = Tail;
};

template <typename List>
struct list_length;

template <>
struct list_length<empty> {
  static const size_t value = 0;
};

template <typename Head, typename Tail>
struct list_length<append<Head, Tail>> {
  static const size_t value = 1 + list_length<Tail>::value;
};
\end{tcblisting}
\end{itemize}
\end{frame}

\begin{frame}[fragile]
\frametitle{Compile-time вычисления: constexpr}
\begin{itemize}
\usemintedstyle{tango}
\item Ключевое слово \mintinline{cpp}|constexpr| означает, что переменная или функция \textit{могут} быть вычислены в compile-time
\pause
\usemintedstyle{lightbulb}
\begin{tcblisting}{listing engine=minted, minted language=cpp, minted options={fontsize=\scriptsize}, listing only, colback=codebg, colframe=codebg, rounded corners}
constexpr unsigned int factorial(unsigned int n) {
  if (n == 0)
    return 1;
  return n * factorial(n - 1);
}
\end{tcblisting}
\pause
\item На такую функцию есть ограничения (не может быть виртуальной, не может иметь побочные эффекты, и т.п.)
\end{itemize}
\end{frame}

\begin{frame}[fragile]
\frametitle{const, constexpr, consteval, constinit}
\begin{itemize}
\usemintedstyle{tango}
\item \mintinline{cpp}|constexpr| переменная будет автоматически \mintinline{cpp}|const|
\pause
\item \mintinline{cpp}|consteval| функция \textbf{всегда} вычисляется в compile-time
\pause
\item \mintinline{cpp}|constinit| переменная инициализируется в compile-time (но не будет автоматически \mintinline{cpp}|const|)
\end{itemize}
\end{frame}

\begin{frame}[fragile]
\frametitle{Variadic templates}
\begin{itemize}
\usemintedstyle{tango}
\item Шаблон (функции или класса), параметризующийся \textit{произвольным списком типов (произвольной длины)}
\pause
\item Синтаксис: \mintinline{cpp}|<typename ... Args>|
\pause
\usemintedstyle{lightbulb}
\begin{tcblisting}{listing engine=minted, minted language=cpp, minted options={fontsize=\scriptsize}, listing only, colback=codebg, colframe=codebg, rounded corners}
template <typename ... Args>
void print(Args ... args);
\end{tcblisting}
\pause
\usemintedstyle{tango}
\item К типам можно добавить уточнения (ссылки, \mintinline{cpp}|const|):
\usemintedstyle{lightbulb}
\begin{tcblisting}{listing engine=minted, minted language=cpp, minted options={fontsize=\scriptsize}, listing only, colback=codebg, colframe=codebg, rounded corners}
template <typename ... Args>
void print(const Args& ... args);
\end{tcblisting}
\end{itemize}
\end{frame}

\begin{frame}[fragile]
\frametitle{Variadic templates}
\begin{itemize}
\usemintedstyle{tango}
\item Сам набор аргументов \mintinline{cpp}|args| называется \textit{variadic pack}
\pause
\item Выражение \mintinline{cpp}|args...| позволяет `раскрыть' variadic pack и передать куда-нибудь (напр. в функцию)
\pause
\item \mintinline{cpp}|sizeof...(Args)| -- количество элементов в наборе
\end{itemize}
\end{frame}

\begin{frame}[fragile]
\frametitle{Variadic templates}
\begin{itemize}
\usemintedstyle{tango}
\item Два способа использовать variadic pack: рекурсия и fold expressions
\pause
\item Пример рекурсии:
\usemintedstyle{lightbulb}
\begin{tcblisting}{listing engine=minted, minted language=cpp, minted options={fontsize=\scriptsize}, listing only, colback=codebg, colframe=codebg, rounded corners}
// Типичный подход: первый аргумент отдельно,
// рекурсия по оставшимся аргументам
template <typename T, typename ... Args>
void print(const T& first, const Args& ... args) {
  std::cout << first;
  print(args...);
}

// База рекурсии
void print() {}
\end{tcblisting}
\end{itemize}
\end{frame}

\begin{frame}[fragile]
\frametitle{Fold expressions}
\begin{itemize}
\usemintedstyle{tango}
\item Правая свёртка: выражение \mintinline{cpp}|(args + ...)| превращается в \mintinline{cpp}|arg0 + (arg1 + (arg2 + ...))|
\pause
\item Левая свёртка: выражение \mintinline{cpp}|(... + args)| превращается в \mintinline{cpp}|((arg0 + arg1) + arg2) + ...|
\pause
\item Правая свёртка с начальным значением: выражение \mintinline{cpp}|(args + ... + init)| превращается в \mintinline{cpp}|arg0 + (arg1 + (arg2 + ... + (argn + init)))|
\pause
\item Левая свёртка с начальным значением: выражение \mintinline{cpp}|(init ... + args)| превращается в \mintinline{cpp}|(((init + arg0) + arg1) + arg2) + ...|
\end{itemize}
\end{frame}

\begin{frame}[fragile]
\frametitle{Fold expressions}
\begin{itemize}
\usemintedstyle{lightbulb}
\begin{tcblisting}{listing engine=minted, minted language=cpp, minted options={fontsize=\scriptsize}, listing only, colback=codebg, colframe=codebg, rounded corners}
template <typename ... Args>
auto sum(const Args& ... args) {
  return (args + ...);
}
\end{tcblisting}
\pause
\begin{tcblisting}{listing engine=minted, minted language=cpp, minted options={fontsize=\scriptsize}, listing only, colback=codebg, colframe=codebg, rounded corners}
template <typename ... Args>
void print(const Args& ... args) {
  (std::cout << ... << args);
}
\end{tcblisting}
\end{itemize}
\end{frame}

\begin{frame}[fragile]
\frametitle{Fold expressions}
\begin{itemize}
\usemintedstyle{tango}
\item Свёртку можно использовать почти со всеми операторами, включая оператор запятая -- это позволяет реализовать `цикл по variadic pack':
\usemintedstyle{lightbulb}
\begin{tcblisting}{listing engine=minted, minted language=cpp, minted options={fontsize=\scriptsize}, listing only, colback=codebg, colframe=codebg, rounded corners}
template <typename ... Args>
void print(const Args& ... args) {
  (std::cout << args, ...);
}
\end{tcblisting}
\begin{tcblisting}{listing engine=minted, minted language=cpp, minted options={fontsize=\scriptsize}, listing only, colback=codebg, colframe=codebg, rounded corners}
template <typename ... Args>
void clear_all(const Args& ... args) {
  (args.clear(), ...);
}
\end{tcblisting}
\end{itemize}
\end{frame}

\end{document}